\section{The Second-Moment Correction: Duan Smearing}
\label{sec:smearing}

% Arc: developed HERE, immediately after the data section introduces it,
% because every downstream fit comparison (buckets, 2-block, 3-block, 4-block)
% inherits whatever smear machinery this section lands on. Raw-space QLIKE
% depends on the sqrt-space -> raw-space back-transform; the second-moment
% correction IS part of the evaluation contract, so it must be frozen before
% the first model comparison.
%
% Planned upgrade, same slot: if/when an LSTM is fit to the correction
% (a learned second-moment model replacing the empirical smear), it wedges in
% right here — between data prep and the bucket attribution — so the contract
% upgrade happens before any comparison that would inherit it.
% Source material: reference/sections/methodology.tex (Duan smearing),
% src/evaluation/metrics.py.

\subsection{Why the Back-Transform Needs a Second Moment}
\label{sec:smearing_bias}

Every model in this paper is fit on the preprocessed scale: realized variance
is diurnally adjusted and then passed through the square-root transform, so
the regression target is $y_t = \sqrt{\widetilde{RV}_t}$. Evaluation, however,
happens on the raw variance scale, because QLIKE --- our primary loss --- is a
strictly consistent scoring rule for the conditional \emph{variance}, not for
any transform of it. A forecast issued in square-root space must therefore be
mapped back to raw space before it can be scored, and it is at this seam that
a subtlety enters.

The na\"ive inverse is to square the forecast and undo the diurnal
adjustment. But squaring does not commute with conditional expectation: by
Jensen's inequality, $\mathbb{E}[\sqrt{X}]^2 \le \mathbb{E}[X]$, with
equality only in the degenerate case. Concretely, if we decompose the target
as $y_{t} = \hat{y}_{t} + \varepsilon_{t}$ with a conditionally mean-zero
innovation, then
\begin{equation}
    \mathbb{E}\!\left[y_t^2 \,\middle|\, \mathcal{F}\right]
    = \hat{y}_t^2 + \mathbb{E}\!\left[\varepsilon_t^2 \,\middle|\, \mathcal{F}\right],
    \label{eq:second_moment}
\end{equation}
so the conditional mean of the squared target exceeds the squared forecast by
exactly the second moment of the square-root-scale innovation. A forecast that
is well calibrated in $\sqrt{RV}$-space is, once squared, systematically
biased low in $RV$-space, and the size of the bias is not incidental --- it is
the residual variance of the very model being evaluated.

For QLIKE this downward bias is doubly consequential, because the loss
penalizes under-prediction more heavily than over-prediction. An uncorrected
squared forecast would not merely carry a bias; it would carry a bias in the
direction the loss punishes hardest, and by an amount that differs across
models with different residual variances. The back-transform therefore needs
an explicit second-moment term, and how that term is estimated is part of the
evaluation itself.

\subsection{The Empirical Smear}
\label{sec:smearing_estimator}

We supply the missing second moment with Duan's smearing estimator
\citep{Duan1983}, in the analytic form implemented in
\texttt{src/evaluation/metrics.py}. Let $\{(\hat{y}_t, y_t)\}_{t=1}^{N}$ be
the forecasts and realizations on the transformed (square-root, diurnally
adjusted) scale over an evaluation window, and let $B_t$ denote the diurnal
baseline factor that reverses the Stage-1 adjustment at bar $t$. The
implementation computes a single scalar smear from the window's own forecast
residuals,
\begin{equation}
    \hat{\sigma}^2_{\varepsilon}
    \;=\; \frac{1}{N}\sum_{t=1}^{N}\left(y_t - \hat{y}_t\right)^2,
    \label{eq:smear_scalar}
\end{equation}
and maps both forecast and realization to the raw scale as
\begin{equation}
    \widehat{RV}_t^{\mathrm{raw}}
    = \left(\hat{y}_t^{\,2} + \hat{\sigma}^2_{\varepsilon}\right) B_t,
    \qquad
    RV_t^{\mathrm{raw}} = y_t^{\,2}\, B_t.
    \label{eq:smear_map}
\end{equation}
Equation~\eqref{eq:smear_map} is the analytic counterpart of Duan's empirical
smearing average $\frac{1}{N}\sum_i (\hat{y}_t + \hat{\varepsilon}_i)^2$; the
two coincide when the residuals are mean-zero, since
$\mathbb{E}[(\hat{y}+\varepsilon)^2] = \hat{y}^2 +
2\hat{y}\,\mathbb{E}[\varepsilon] + \sigma^2_{\varepsilon}$. Because the
correction is non-negative, the smeared forecast is positive whenever the
baseline is, so the QLIKE ratio in the next sections is always well defined.

Three implementation choices deserve to be stated plainly, because they are
choices and not forced moves.

\emph{Global, not rolling.} The smear is one scalar per evaluation window,
not a rolling or bar-local estimate. In the chunked walk-forward harness the
backtest is split into contiguous chunks and
Equation~\eqref{eq:smear_scalar} is computed once per chunk, over all of that
chunk's out-of-sample bars; the correction is therefore piecewise-constant
across the panel on the transformed scale. On the raw scale it still varies
bar by bar through the baseline factor $B_t$, and in relative terms it is
largest exactly where the squared forecast is smallest --- the quiet bars.

\emph{Out-of-sample residuals of the model itself.} The residuals entering
Equation~\eqref{eq:smear_scalar} are the evaluated model's own out-of-sample
forecast errors on the window being converted --- not training-set residuals,
and not the residuals of any reference model. Two consequences follow. First,
the smear is model-specific: an arm with larger residual variance receives a
larger upward correction, so the correction is properly understood as part of
each arm's forecast rather than a shared post-processing constant. Second,
the smear at bar $t$ uses residuals from the whole window, including bars
after $t$. This is an evaluation-time convention, applied identically to every
arm, not a quantity a deployed forecaster could have computed at $t$; a
deployable variant would estimate $\hat{\sigma}^2_{\varepsilon}$ from trailing
residuals only, and we flag the distinction here because the correction sits
inside every raw-scale number that follows.

\emph{Two biases left uncorrected.} The residuals in
Equation~\eqref{eq:smear_scalar} are computed on the winsorized scale while
the evaluation target is unwinsorized, and the Stage-3 winsorization is a
lossy operation that is not inverted at all. We do not quantify either
effect; both are shared by every arm and so cancel to first order in
comparisons, but they mean the raw-scale forecasts should be read as
evaluation objects, not as unbiased variance estimates in an absolute sense.

\subsection{What Downstream Inherits}
\label{sec:smearing_contract}

The reason this section precedes any model comparison is that
Equation~\eqref{eq:smear_map} is not a detail of one table --- it is the
common gateway through which every forecast in this paper passes on its way
to a QLIKE number. Every downstream comparison
\pending{enumerate the downstream analyses by their final section labels
(bucket attribution, block structure, model classes) once those sections
stabilize} scores its arms on raw-scale pairs produced by exactly this map,
with exactly these window and residual choices. Whatever the smear does, it
does to all of them; whatever it does differentially, it does through each
arm's own residual variance.

That differential channel is the substantive point. Because
$\hat{\sigma}^2_{\varepsilon}$ is model-specific, the smear does not shift all
arms by a common constant: it inflates weaker models' forecasts more than
stronger ones', and under QLIKE's asymmetry this interacts with each arm's
calibration in a way that a shared correction would not. A comparison run
under a different convention --- training-set residuals, a rolling window, or
no correction at all --- would be a comparison of different forecast objects,
and there is no guarantee it would induce the same ranking.
\pending{sensitivity check: re-score the frozen battery under alternative
smear conventions (no smear, trailing-window smear, training-residual smear)
and report whether any pairwise ranking flips; this is the empirical teeth of
the contract claim}

We therefore treat the correction as part of the evaluation contract: it is
frozen here, before the first comparison, and no downstream section revisits
it. Any future change to the second-moment model --- including the learned
correction contemplated for this slot --- would constitute a change of
contract, and would require re-scoring every comparison that inherited the
old one.
